\subsection{Examples in $\SL_{2}$}
By Theorem \ref{th:orb coho}, in order to calculate orbital integrals, we need to know not just the geometry of the affine Springer fiber $\sX_{\gamma}$, but also the action of Frobenius on its cohomology. Having already seen many examples of affine Springer fibers over an algebraically closed field in \S\ref{ss:aff examples}, our emphasis here will be on the Frobenius action. 

In this subsection we let $G=\SL_{2}$ and assume $\chark>2$. We will compute several orbital integrals in this case and verify Theorem \ref{th:orb coho} in these cases by explicit calculations.

\subsubsection{Unramified case: $\gamma$}\label{unram sl2} Let $\gamma=\mat{0}{at}{t}{0}\in\sl_{2}(F)$ be a regular semisimple element with $a\in k^{\times}-(k^{\times})^{2}$. 

Let $E\subset\mathfrak{gl}_{2}(F)$ be the centralizer of $\gamma$ in $\mathfrak{gl}_{2}(F)$. Then $E=\{u+v\gamma|u,v\in F\}$ is an unramified quadratic extension of $F$ obtained by adjoining $\sqrt{a}$. Therefore we have $E=k_{E}\lr{t}$, with $k_{E}= k(\sqrt{a})$. We have $G_{\gamma}(F)=(E^{\times})^{\Nm=1}=\ker(\Nm: E^{\times}\to F^{\times})$, which is compact. We fix a Haar measure on $G_{\gamma}(F)$ with total volume 1. 

Let $\sX_{\gamma}$ be the affine Springer fiber of $\gamma$, which is a scheme over $k$. Lemma \ref{l:Ogamma} implies that
\begin{equation*}
O_{\gamma}(\one_{\frg(\calO_{F})})=\#X_{\gamma}=\#\sX_{\gamma}(k).
\end{equation*}
In \S\ref{P1 chain} we have shown that $\sX_{\gamma,\kbar}$ is an infinite union of rational curves $C_{n}\cong\PP^{1}$ indexed by the integers $n\in\ZZ$. Since $\gamma$ is diagonalizable over $E$,  each component $C_{n}$ is in fact defined over $k_{E}=k(\sqrt{a})$. The lattice $\L_{\gamma}\subset G_{\gamma}(F^{\ur})$ is contained in $G_{\gamma}(E)\cong E^{\times}$, and is generated by the uniformizer $t\in E$. We label the components $C_{n}$ so that $t\in \L_{\gamma}$ sends $C_{n}$ to $C_{n+1}$. Let $x_{n+1/2}:=C_{n}\cap C_{n+1}$, which is a $k_{E}$-point of $\sX_{\gamma}$.  

The action of  the nontrivial involution $\sigma\in\Gal(k_{E}/k)$ on $G_{\gamma}(E)$ is by inversion, hence it also acts on $\L_{\gamma}$ by inversion.  The standard lattice $\calO^{2}_{F}$ lies in both $C_{0}$ and $C_{1}$, hence it is the point $x_{1/2}$. Therefore the point $x_{1/2}$ is fixed by $\sigma$ since it is defined over $k$. Since the action of $\sigma$ on $\sX_{\gamma,k_{E}}$ is compatible with its action on $\L_{\gamma}$ (by inversion), the only possibility is that
\begin{equation*}
\sigma(C_{n})=C_{1-n}, \quad \sigma(x_{n+1/2})=x_{-n+1/2}, \quad\forall n\in\ZZ.
\end{equation*}
The action of $\sigma$ can be represented by the picture
\begin{equation*}
\xymatrix{ x_{-3/2} \ar@/^{1pc}/[rrrr]_{\sigma}
\ar@{=} '[dr]_{C_{-1}} '[rr]_{C_{0}}  '[drrr]_{C_{1}}  [rrrr]_{C_{2}}
&  & x_{1/2}  & & x_{5/2} \\
& x_{-1/2}\ar@/_1pc/[rr]^{\sigma} &  & x_{3/2}}
\end{equation*}
Here each double line represents a $\PP^{1}$. From this we see that $X_{\gamma}=\sX_{\gamma}(k)$ consists of only one point $x_{1/2}$, namely the standard lattice $\calO^{2}_{F}$. This implies that
\begin{equation}\label{Oga}
O_{\gamma}(\one_{\frg(\calO_{F})})=1.
\end{equation}


\subsubsection{Unramified case: $\gamma'$} Now consider the element $\gamma'=\mat{0}{at^{2}}{1}{0}\in\sl_{2}(F)$. In Exercise \ref{exIII:stable conj} we see that $\gamma'$ is stably conjugate to $\gamma$ but not conjugate to $\gamma$ under $\SL_{2}(F)$. However, $\gamma'$ is conjugate to $\gamma$ under $\SL_{2}(E)$. Therefore, the affine Springer fiber $\sX_{\gamma'}$ still looks the same as $\sX_{\gamma}$ over $k_{E}$, but the action of $\sigma\in\Gal(k_{E}/k)$ is different. 


Consider the component of $\sX_{\gamma',k_{E}}$ whose $k_{E}$-points consist of $\gamma'$-stable lattices $\L\subset E^{2}$ such that $t\calO_{E}\oplus\calO_{E}\subset \L\subset \calO_{E}\oplus t^{-1}\calO_{E}$. This component is cut out by conditions defined over $k$, so it is stable under $\sigma$, and we call this component $C'_{0}$. We label the other components of $\sX_{\gamma',k_{E}}$ by $C'_{n}$ ($n\in\ZZ$) so that the generator $t\in\L_{\gamma'}$ sends $C'_{n}$ to $C'_{n+1}$. Let $x'_{n+1/2}=C'_{n}\cap C'_{n+1}\in \sX_{\gamma'}(k_{E})$. 
Since the action of $\sigma$ on $\sX_{\gamma,k_{E}}$ is compatible with its action on $\L_{\gamma'}$ (by inversion), the only possibility is that
\begin{equation*}
\sigma(C'_{n})=C'_{-n}, \quad \sigma(x'_{n+1/2})=x'_{-n-1/2},\quad\forall n\in\ZZ.
\end{equation*}
The action of $\sigma$ can be represented by the picture
\begin{equation*}
\xymatrix{ & & x'_{-1/2} \ar@/_{1pc}/@<-1ex>[r]_{\sigma}& x'_{1/2}\\
x'_{-5/2}  \ar@/_2pc/[rrrrr]^{\sigma} 
\ar@{=} '[r]^{C'_{-2}} '[urr]_{C'_{-1}}  '[urrr]_{C'_{0}}  '[rrrr]_{C'_{1}} [rrrrr]^{C'_{2}}
&  x'_{-3/2} \ar@/_{1pc}/[rrr]^{\sigma}& & & x'_{3/2}  & x'_{5/2}}
\end{equation*}
Therefore no point $x'_{n+1/2}$ is defined over $k$. The component $C'_{0}$ is the only one that is defined over $k$, and it has to be isomorphic to $\PP^{1}$ over $k$ because it is so over $k_{E}$ (there are no nontrivial Brauer-Severi varieties over a finite field). We see that $\sX_{\gamma'}(k)=C'_{0}(k)\cong\PP^{1}(k)$ has $q+1$ elements. Therefore
\begin{equation}\label{Ogb}
O_{\gamma'}(\one_{\frg(\calO_{F})})=q+1.
\end{equation}
Adding up \eqref{Oga} and \eqref{Ogb} we get
\begin{equation*}
SO_{\gamma}(\one_{\frg(\calO_{F})})=q+2.
\end{equation*}
 
\subsubsection{Unramified case: cohomology} The quotient $\L_{\gamma}\bs\sX_{\gamma}$ is a nodal rational curve obtained from $\PP^{1}$ by glueing two $k$-points into a nodal point. 

Now let us consider the quotient $\L_{\gamma'}\bs\sX_{\gamma'}$. Over $k_{E}$ this is also a nodal rational curve consisting of a unique node $y$ which is image of all $x_{j}$. While $y$ is a $k$-point of the quotient $\L_{\gamma'}\bs\sX_{\gamma'}$, none of its preimages $x_{j}$ are defined over $k$. Therefore, the $q+1$ points in $C'_{0}(k)$ still map injectively to the quotient, in addition to the point $y$. We conclude that $(\L_{\gamma'}\bs\sX_{\gamma'})(k)$ consists of $q+2$ points.  Since $LG_{\gamma'}/\L_{\gamma'}$ is connected, the stable part of the cohomology of $\L_{\gamma'}\bs\sX_{\gamma'}$ is the whole $\cohog{*}{\L_{\gamma'}\bs\sX_{\gamma'}}$, and the alternating sum of Frobenius trace on it is the cardinality of $(\L_{\gamma'}\bs\sX_{\gamma'})(k)$. In this special case we have verified the formula \eqref{d}. We remark that the action of $\sigma$ on the 1-dimensional space $\cohog{1}{\L_{\gamma'}\bs\sX_{\gamma',\kbar},\Ql}$ is by $-1$, and the Grothendieck-Lefschetz trace formula for Frobenius reads $\#(\L_{\gamma'}\bs\sX_{\gamma'})(k)= 1+1+q$ instead of $1-1+q=q$, the latter being the number of $k$-points on a nodal rational curve obtained by identifying two $k$-points on $\PP^{1}$.


\subsubsection{Ramified case: orbital integrals}\label{sss:ram orb SL2} Consider the elements $\gamma=\mat{0}{t^{2}}{t}{0}$ and $\gamma'=\mat{0}{at^{2}}{a^{-1}t}{0}$ where $a\in k^{\times}-(k^{\times})^{2}$. Again $\gamma$ and $\gamma'$ are stably conjugate but not conjugate in $\SL_{2}(F)$. 

Let $E$ be the centralizer of $\gamma$ in $\gl_{2}(F)$. Then $E$ is a ramified quadratic extension of $F$, and $G_{\gamma}(F)=(E^{\times})^{\Nm=1}$. Similarly, let $E'$ be the centralizer of $\gamma'$ in $\gl_{2}(F)$. Then $E'$ is another ramified quadratic extension of $F$, and $G_{\gamma'}(F)=(E'^{\times})^{\Nm=1}$. We choose Haar measures on compact groups $G_{\gamma}(F)$ and $G_{\gamma'}(F)$ with total volume $1$.  

In both cases, $\sX_{\gamma}$ and $\sX_{\gamma'}$ are isomorphic to $\PP^{1}$ as varieties over $k$. In fact we have shown in \S\ref{SL2 ram dim 1} that these varieties are isomorphic to $\PP^{1}$ over $\kbar$, hence they must be isomorphic to $\PP^{1}$ over $k$ as well.  Both $\sX_{\gamma}(k)$ and $\sX_{\gamma'}(k)$  consist of lattices $t\calO_{F}\oplus\calO_{F}\subset \L\subset \calO_{F}\oplus t^{-1}\calO_{F}$, therefore $\sX_{\gamma}=\sX_{\gamma'}$ as subvarieties of $\Gr_{G}$. By Lemma \ref{l:Ogamma}, we have
\begin{equation*}
O_{\gamma}(\one_{\frg(\calO_{F})})=O_{\gamma'}(\one_{\frg(\calO_{F})})=\#\sX_{\gamma}(k)=q+1.
\end{equation*}
Therefore
\begin{equation}\label{ram so}
SO_{\gamma}(\one_{\frg(\calO_{F})})=2(q+1).
\end{equation}


\subsubsection{Ramified case: cohomology} In the setup of \S\ref{sss:ram orb SL2}, $\L_{\gamma}=0$. The component group of $LG_{\gamma}$ is $\ZZ/2\ZZ$, but its action on $\sX_{\gamma}$ is trivial. Therefore the stable part of the cohomology is the whole $\cohog{*}{\sX_{\gamma},\Ql}$, on which the alternating sum of the Frobenius gives the cardinality of $\sX_{\gamma}(k)\cong\PP^{1}(k)$.
However, the parahoric subgroup of $G_{\gamma}(F)=(E^{\times})^{\Nm=1}$ has index $2$ in it ($K_{\gamma}$ consists of those $e\in (E^{\times})^{\Nm=1}$ whose reduction in $k$ is $1$). Therefore, the right side of formula \eqref{d} gets a factor $2=\vol(K_{\gamma},d_{\gamma}g)^{-1}$ in front of $\#\sX_{\gamma}(k)$. This is consistent with \eqref{ram so}, and we have checked the formula \eqref{d} in our special case.

\subsection{Remarks on the Fundamental Lemma} Let us go back to the situation in  \S\ref{unram sl2}. What happens if we take the difference of $O_{\gamma'}(\one_{\frg(\calO_{F})})$ and $O_{\gamma}(\one_{\frg(\calO_{F})})$ instead of their sum? Is there a geometric interpretation of this difference analogous to Theorem \ref{th:orb coho}?

\subsubsection{The $\kappa$-orbital integral} The linear combination $O_{\gamma}(\one_{\frg(\calO_{F})})-O_{\gamma'}(\one_{\frg(\calO_{F})})$ is an example of $\kappa$-orbital integrals. More generally, let $\kappa$ be a character of $\cohog{1}{F,G_{\gamma}}$, then we define the $\kappa$-orbital integral of $\varphi\in \sS(\frg(F))$ to be
\begin{equation*}
O^{\kappa}_{\gamma}(\varphi)=\sum_{\gamma'}\kappa(\inv(\gamma,\gamma'))O_{\gamma'}(\varphi)
\end{equation*}
where the sum is over the $G(F)$-orbits in the stable conjugacy class of $\gamma$.


\subsubsection{Statement of the Fundamental Lemma} The Langlands-Shelstad conjecture, also known as the {\em Fundamental Lemma}, states that the $\kappa$-orbital integral of $\one_{\frg(\calO_{F})}$ for $\gamma\in\frg(F)$ is equal to the stable orbital integral of an element $\gamma_{H}\in\frh(F)$ for a smaller group $H$, up to a simple factor. In formula, the Fundamental Lemma is the identity
\begin{equation*}
O^{\kappa}_{\gamma}(\one_{\frg(\calO_{F})})=\Delta(\gamma,\gamma_{H})\cdot SO_{\gamma_{H}}(\one_{\frh(\calO_{F})}).
\end{equation*}
The smaller group $H$ depends on both $G$ and $\kappa$, and is called the {\em endoscopic group} of $(G,\kappa)$. The number $\Delta(\gamma,\gamma_{H})$ is called the {\em transfer factor}, which turns out to be an integer power of $q$ (depending on $\gamma$ and $\gamma_{H}$) if $\gamma$ is chosen appropriately from its stable conjugacy class, and the measures on $G_{\gamma}$ and $H_{\gamma_{H}}$ are chosen properly.


\subsubsection{A simple case} In the situation of \S\ref{unram sl2}, take the nontrivial character $\kappa$ on $\cohog{1}{F,G_{\gamma}}\cong\ZZ/2\ZZ$, the corresponding endoscopic group $H$ is isomorphic to the torus $G_{\gamma}$; but in general it is not always isomorphic to a subgroup of $G$. The Fundamental Lemma in this case is the identity
\begin{equation*}
O_{\gamma'}(\one_{\frg(\calO_{F})})-O_{\gamma}(\one_{\frg(\calO_{F})})=q=q\cdot SO_{\gamma_{H}}(\one_{\frh(\calO_{F})})
\end{equation*}
where $\gamma_{H}=\gamma$ if we identify $H$ with $G_{\gamma}$. 

On the other hand, in the ramified situation \S\ref{sss:ram orb SL2}, the $\kappa$-orbital integral of $\gamma$ for the nontrivial $\kappa$ vanishes. This maybe explained without calculating the orbital integrals explicitly, for in general, $\kappa$ must factor through a further quotient of $\cohog{1}{F,G_{\gamma}}$ for $O^{\kappa}_{\gamma}(\one_{\frg(\calO_{F})})$ to be possibly nonzero. For the precise statement, see \cite[Prop 8.2.7]{NgoFL}.

\subsubsection{Comments on the proof} The Fundamental Lemma for general $G$ and function field $F$ was established by B-C. Ng\^o \cite{NgoFL}. There is a generalization of Theorem \ref{th:orb coho} to $\kappa$-orbital integrals, in which we replace the stable part of the cohomology of $\L_{\gamma}\bs\sX_{\gamma}$ by the $\kappa$-isotypic part. Using this generalization, Goresky, Kottwitz and MacPherson \cite{GKM} reformulated the Fundamental Lemma as a relation between cohomology groups of the affine Springer fiber $\sX_{\gamma}$ and its endoscopic cousin $\sX^{H}_{\gamma_{H}}$. They were also able to prove the Fundamental Lemma in some special but highly nontrivial cases. Ng\^o's proof builds on this cohomological reformulation, but also uses a new ingredient, namely Hitchin fibers, which can be viewed as a ``global'' analog of affine Springer fibers. This will be the topic of the next lecture.  


\subsection{Exercises}

In these exercises, $k=\FF_{q}$ denotes a finite field with $\chark\neq2$, and $F=k\lr{t}$.


\subsubsection{}\label{exIII:K inv} Let $H$ be an algebraic group over $F$ and let $f\in\sS(H(F))$. Show that there exists a compact open subgroup $K\subset H(F)$ such that $f$ is both left and right invariant under $K$.

\subsubsection{}\label{exIII:Ogamma} Prove Lemma \ref{l:Ogamma}.


\subsubsection{}\label{exIII:stable conj} Let $a\in k^{\times}-(k^{\times})^{2}$ and $G=\SL_{2}$. Consider the matrices $\gamma,\gamma'$ as in \S\ref{sss:st conj SL2}. Show that  $\gamma$ and $\gamma'$ are stably conjugate but not conjugate under $\SL_{2}(F)$.


\subsubsection{} Let $G=\GL_{n}$ over $F$ and let $\gamma=\diag(\gamma_{1},\cdots,\gamma_{m})$ be a block diagonal matrix in $\frg(F)^{\rs}$, where $\gamma_{i}\in\mathfrak{gl}_{n_{i}}(F)^{\rs}$. Consider the asymptotic behavior of the orbital integral $O_{\gamma}(\one_{\frg(\calO_{F})})$  as $q=\#k$ tends to $\infty$. Find the smallest integer $d$ such that
\begin{equation*}
O_{\gamma}(\one_{\frg(\calO_{F})})=O(q^{d})
\end{equation*}
as $q\to\infty$. Note that the $O(\cdot)$ on the right side is analysts' $O$ while on the left side it means orbital integral.  Can you interpret $d$ in terms of the characteristic polynomials of the $\gamma_{i}$'s? 

For an explicit estimate of $O_{\gamma}(\one_{\frg(\calO_{F})})$, see \cite{Y-Orb}.



\subsubsection{} Let $G=\SL_{3}$ and $\gamma=\left(\begin{array}{ccc} 0 & 0 & t^{4}\\ 1 & 0 & 0 \\ 0 & 1& 0 \end{array}\right)\in\frg(F)$. Compute $O_{\gamma}(\one_{\frg(\calO_{F})})$. 

{\em Hint}: use the cell decomposition introduced in \S\ref{sss:LS}.
 
 
\subsubsection{} Let $G=\SL_{2}$. Let $f$ be the characteristic function of elements $X\in\frg(\calO_{F})$ such that the reduction $\ov{X}$ in $\frg(k)$ is regular nilpotent. Let $\gamma\in\frg(F)$ be a regular semisimple element.  
\begin{enumerate}
\item Show that $O_{\gamma}(f)=0$ unless $\det(\gamma)\in t\calO_{F}$.
\item When $\det(\gamma)\in t\calO_{F}$, show that
\begin{equation*}
O_{\gamma}(f)=O_{\gamma}(\one_{\frg(\calO_{F})})-O_{t^{-1}\gamma}(\one_{\frg(\calO_{F})}).
\end{equation*}
\end{enumerate}

\subsubsection{} Let $G=\GL_{2}$ and $\gamma=\mat{0}{t^{n}}{1}{0}$ for $n\geq 1$ odd. Let $G(F)_{d}$ be the set of $g\in G(F)$ with $\val_{F}(\det g)=d$. Fix the Haar measure on $G(F)$ such that $G(\calO_{F})$ has volume $1$. Show that, for any integer $d\geq0$,
\begin{equation*}
\int_{G(F)_{d}}\one_{\frg(\calO_{F})}(g^{-1}\gamma g)\one_{\calO^{2}_{F}}((0,1)g)dg
\end{equation*}
is the same as the number of closed subschemes $Z$ of the plane curve $y^{2}-t^{n}=0$ satisfying: (1) the underlying topological space of $Z$ is the point $(y,t)=(0,0)$; (2) $\dim_{k}\calO_{Z}=d$.

Note: this exercise relates orbital integrals to Hilbert schemes of points on curves. This relationship has been used in \cite{Y-Orb} to provide an estimate for orbital integrals for $\GL_{n}$. See also \S\ref{sss:MHilb} and \S\ref{sss:Hilb} for an global analog.

%%%%%%%%%%%%%%%%%%%%%%


